\begin{proposition}[2.4.4]
\label{IV.2.4.4-fr}
Un morfismo universalmente bicontinuo $f:X\to Y$ es radicial (y por tanto
separado \textup{(1.8.7.1)}).
\end{proposition}

En efecto, $f$ es universalmente inyectivo por hipótesis (\textbf I, 3.5.11).

\begin{proposition}[2.4.5]
\EGAReaderTarget{ega4.c04.prop.2-4-5}{IV.2.4.5-fr}
\begin{enumerate}
\item[\textup{(i)}] Un morfismo $f:X\to Y$ que es entero, sobreyectivo y radicial es un
homeomorfismo universal.
\item[\textup{(ii)}] Recíprocamente, supongamos que $Y$ es localmente noetheriano.
Entonces, si un morfismo de tipo finito $f:X\to Y$ es un homeomorfismo
universal, es finito, sobreyectivo y radicial.
\end{enumerate}
\end{proposition}

(i) Basta observar que las tres propiedades de $f$ se conservan por cambio de
base (\textbf I, 3.5.2, \textbf{I}, 3.5.7 y \textbf{II}, 6.1.5), y, como un morfismo entero es cerrado
(\textbf{II}, 6.1.10), es claro que $f$ es un homeomorfismo de $X$ sobre $Y$.

(ii) Como $f$ es de tipo finito y universalmente cerrado por hipótesis y es
separado por (2.4.4), es propio (\textbf{II}, 5.4.1), y, para todo $y\in Y$,
$f^{-1}(y)$ se reduce a un único elemento; por tanto (\textbf{III}, 4.4.2), $f$ es
finito; es claro que $f$ es sobreyectivo, y es radicial puesto que es
universalmente inyectivo (\textbf I, 3.5.11).

\begin{theorem}[2.4.6]
\label{IV.2.4.6-fr}
Sea $f:X\to Y$ un morfismo cuasiplano \textup{(2.3.3)} y localmente de presentación
finita. Entonces $f$ es universalmente abierto. En particular, un morfismo
plano localmente de presentación finita es universalmente abierto.
\end{theorem}

Se sabe que, para todo cambio de base $Y'\to Y$, $f_{(Y')}$ es cuasiplano
(2.3.3) y localmente de presentación finita (1.4.3, (iii)). Basta, pues,
demostrar que $f$ es un
\oldpage[IV-2]{21}
morfismo \emph{abierto}. Pero esto se deduce del criterio (1.10.4) para los morfismos
localmente de presentación finita, pues las condiciones \emph{b)},
\emph{b$'$)} y \emph{c)} de (1.10.4) no son otras que las condiciones (i),
(ii), (iii) de (2.3.4).

\begin{corollary}[2.4.7]
\label{IV.2.4.7-fr}
Para todo preesquema $Y$, el morfismo estructural $\mathbf V_Y^n\to Y$ \textup{(donde}
$\mathbf V_Y^n
=Y\otimes_{\mathbf Z}\operatorname{Spec}
  \bigl(\mathbf Z[T_1,\ldots,T_n]\bigr)$,
\textup{también denotado $Y[T_1,\ldots,T_n]$)} es universalmente abierto.
\end{corollary}

En efecto, para $Y=\operatorname{Spec}(A)$,
$Y[T_1,\ldots,T_n]=\operatorname{Spec}(A[T_1,\ldots,T_n])$, y
$A[T_1,\ldots,T_n]$ es una $A$-álgebra libre de presentación finita.

\begin{remarks}[2.4.8]
\label{IV.2.4.8-fr}
\begin{enumerate}
\item[(i)] Se observará que un morfismo $f:X\to Y$ fielmente plano y
cuasicompacto \emph{no es necesariamente abierto}, incluso cuando $X$ e $Y$
son noetherianos. Tomemos, por ejemplo, como $Y$ una curva algebraica reducida
irreducible (\textbf{II}, 7.4) cuyo punto genérico es $y$, y sea $X$ el preesquema suma
$Y\amalg\operatorname{Spec}(\boldsymbol{k}(y))$, siendo $f$ el morfismo canónico; es claro
que $f$ es plano y sobreyectivo, por tanto fielmente plano, y cuasicompacto,
pero la imagen por $f$ de la parte abierta $\operatorname{Spec}(\boldsymbol{k}(y))$ de $X$
es el conjunto $\{y\}$, que no es abierto en $Y$ (\textbf{II}, 7.4.3).
\item[(ii)] Para todo preesquema $X$, el morfismo canónico
$f:X_{\mathrm{red}}\to X$ es una inmersión cerrada y un
\emph{homeomorfismo universal} (2.4.3, (vi)); pero, cuando $X$ es localmente
noetheriano, $f$ solo es \emph{plano} si $X$ es reducido, y por tanto
$f=1_X$ (2.2.17).
\end{enumerate}
\end{remarks}
